\documentclass[preview]{standalone}

\usepackage{amsmath}
\usepackage{amssymb}
\usepackage{stellar}
\usepackage{definitions}

\begin{document}

\id{finitely-generated-modules-over-principal-ideal-domains}
\genpage

\section{Cyclic decomposition}

\begin{snippettheorem}{finitely-generated-pid-module-cyclic-decomposition}{Cyclic decomposition over a PID}
    Let \(A\) be a \principalidealdomain and let \(M\) be a finitely
    generated \(A\)-\module. There exist \(r,s\in\naturalnumbers\) and
    nonzero nonunits \(c_1,\ldots,c_s\in A\) such that
    \[
        M\moduleisomorphic
        A^r\oplus
        A/(c_1)\oplus\cdots\oplus A/(c_s).
    \]
    Moreover,
    \[
        \moduletor(M)
        \moduleisomorphic
        A/(c_1)\oplus\cdots\oplus A/(c_s),
        \qquad
        M/\moduletor(M)\moduleisomorphic A^r.
    \]
\end{snippettheorem}

\begin{snippetproof}{finitely-generated-pid-module-cyclic-decomposition-proof}{finitely-generated-pid-module-cyclic-decomposition}{Cyclic decomposition over a PID}
    Choose generators \(v_1,\ldots,v_n\) of \(M\). The universal
    property of \(A^n\) gives a surjective \modulehomomorphism
    \[
        \varphi\colon A^n\fromto M,
        \qquad
        \varphi(e_i)=v_i.
    \]
    Hence \(M\moduleisomorphic A^n/\moduleker(\varphi)\). By the
    compatible-bases theorem, there are a basis
    \(x_1,\ldots,x_n\) of \(A^n\) and elements
    \(c_1,\ldots,c_n\in A\) such that
    \[
        \moduleker(\varphi)
        =
        \langle c_1x_1,\ldots,c_nx_n\rangle.
    \]
    Therefore
    \[
        M
        \moduleisomorphic
        \bigoplus_{i=1}^{n}A/(c_i).
    \]
    A zero \(c_i\) contributes a copy of \(A\), while a unit \(c_i\)
    contributes the zero module. After discarding the latter and
    collecting the former, the displayed decomposition follows.

    Each quotient by a nonzero \(c_i\) is torsion, whereas \(A^r\) is
    torsion-free. Thus the direct sum of the quotient factors is exactly
    the torsion submodule, and quotienting it out leaves \(A^r\).
\end{snippetproof}

\begin{snippetcorollary}{finitely-generated-torsion-free-pid-module-free}{Finitely generated torsion-free modules over a PID are free}
    Every finitely generated torsion-free \module over a
    \principalidealdomain is free.
\end{snippetcorollary}

\begin{snippetproof}{finitely-generated-torsion-free-pid-module-free-proof}{finitely-generated-torsion-free-pid-module-free}{Finitely generated torsion-free modules over a PID are free}
    In the cyclic decomposition, the sum of all quotient factors is
    \(\moduletor(M)\). If \(M\) is torsion-free, this sum vanishes, so
    \(M\moduleisomorphic A^r\).
\end{snippetproof}

\section{Elementary divisors}

\begin{snippettheorem}{finitely-generated-pid-module-primary-decomposition}{Primary decomposition and elementary divisors}
    Let \(M\) be a finitely generated \module over the
    \principalidealdomain \(A\). Then
    \[
        M\moduleisomorphic
        A^r\oplus
        \bigoplus_{i=1}^{t}A/(p_i^{\alpha_i}),
    \]
    where \(r\in\naturalnumbers\), every \(p_i\) is an
    \irreducibleelement, and every \(\alpha_i\geq1\). Repetitions and
    associates among the \(p_i\) are allowed.

    The integer \(r\) and the multiset of associate classes of the
    powers \(p_i^{\alpha_i}\) are uniquely determined by \(M\), up to
    reordering. These powers are the
    \elementarydivisor[elementary divisors] of
    \(M\). For a fixed \primeelement \(p\), the \primarycomponent \(M_p\)
    is the direct sum of precisely those factors for which \(p_i\) is
    associate to \(p\).
\end{snippettheorem}

\begin{snippetproof}{finitely-generated-pid-module-primary-decomposition-proof}{finitely-generated-pid-module-primary-decomposition}{Primary decomposition and elementary divisors}
    Apply the cyclic decomposition. Factor every \(c_j\) into pairwise
    nonassociate prime powers:
    \[
        c_j=u_jq_{j,1}^{\beta_{j,1}}\cdots
        q_{j,k_j}^{\beta_{j,k_j}}.
    \]
    The ideals generated by distinct prime powers are pairwise comaximal,
    so the Chinese remainder theorem gives
    \[
        A/(c_j)
        \moduleisomorphic
        \bigoplus_{\nu=1}^{k_j}
        A/(q_{j,\nu}^{\beta_{j,\nu}}).
    \]
    This produces the asserted primary decomposition.

    In the resulting direct sum, an element is \(p\)-primary exactly
    when its free coordinates and all coordinates belonging to primes
    nonassociate to \(p\) vanish. Hence \(M_p\) is exactly the sum of
    the \(p\)-power factors. For each fixed \(p\), uniqueness of the
    exponents follows from the successive-socle dimensions. The free
    rank is determined by \(M/\moduletor(M)\). Thus all elementary
    divisors and \(r\) are intrinsic.
\end{snippetproof}

\begin{snippetnote}{primary-decomposition-maximal-cyclic-splitting}{The finest cyclic decomposition}
    A quotient \(A/(p^\alpha)\) cannot be decomposed further as a direct
    sum of nonzero cyclic modules with coprime annihilators. In this
    sense, primary decomposition has the largest possible number of
    cyclic torsion summands.
\end{snippetnote}

\section{Invariant factors}

\begin{snippettheorem}{finitely-generated-pid-module-invariant-factor-decomposition}{Invariant-factor decomposition}
    Let \(M\) be a finitely generated \module over the
    \principalidealdomain \(A\). There exist a unique
    \(r\in\naturalnumbers\) and nonzero nonunits
    \(d_1,\ldots,d_s\in A\) such that
    \[
        d_1\divides d_2\divides\cdots\divides d_s
    \]
    and
    \[
        M\moduleisomorphic
        A^r\oplus A/(d_1)\oplus\cdots\oplus A/(d_s).
    \]
    Each \(d_i\) is unique up to multiplication by a unit. The elements
    \(d_i\) are the \invariantfactor[invariant factors] of \(M\).
\end{snippettheorem}

\begin{snippetproof}{finitely-generated-pid-module-invariant-factor-decomposition-proof}{finitely-generated-pid-module-invariant-factor-decomposition}{Invariant-factor decomposition}
    Start with the elementary-divisor decomposition. For each associate
    class of primes \(p\), list its powers in nondecreasing order of
    exponent. Right-align all these finite lists and multiply the prime
    powers appearing in each column. Call the resulting products
    \(d_1,\ldots,d_s\), from left to right. Exponents are nondecreasing
    along every row, so
    \[
        d_1\divides\cdots\divides d_s.
    \]
    Within each column the prime powers are pairwise coprime. The Chinese
    remainder theorem therefore combines their cyclic quotients into
    \(A/(d_j)\). This proves existence.

    Conversely, factoring every \(d_j\) into prime powers recovers the
    elementary divisors. Their uniqueness therefore implies uniqueness
    of the associate classes of the \(d_j\). The free rank is already
    determined by \(M/\moduletor(M)\).
\end{snippetproof}

\begin{snippetexample}{primary-to-invariant-factor-example}{From elementary divisors to invariant factors}
    Suppose the torsion part of \(M\) has elementary-divisor
    decomposition
    \[
        A/(p_1^3)\oplus A/(p_1^3)\oplus A/(p_1)\oplus A/(p_1)
        \oplus A/(p_2^2)\oplus A/(p_2^2)\oplus A/(p_2)
        \oplus A/(p_3^4),
    \]
    where \(p_1,p_2,p_3\) are pairwise nonassociate primes. Right-aligning
    the powers gives the invariant factors
    \[
        d_1=p_1,
        \qquad
        d_2=p_1p_2,
        \qquad
        d_3=p_1^3p_2^2,
        \qquad
        d_4=p_1^3p_2^2p_3^4.
    \]
    Thus \(d_1\divides d_2\divides d_3\divides d_4\), and the torsion
    part is
    \[
        A/(d_1)\oplus A/(d_2)\oplus A/(d_3)\oplus A/(d_4).
    \]
\end{snippetexample}

\end{document}
